\newpage
\section{Introduzione}

\subsection{Contesto e Scenario Applicativo}
Nell'attuale ecosistema dell'\textit{Internet of Things} (IoT), l'acquisizione sinergica di parametri ambientali e coordinate geografiche rappresenta un requisito abilitante per il monitoraggio remoto avanzato e l'integrazione nei processi di \textit{Industria 4.0}.

Il \textit{data logging} georeferenziato consiste nella registrazione di flussi di dati sensoriali indissolubilmente legati a coordinate spaziali e riferimenti temporali precisi, permettendo la ricostruzione analitica di mappe fenomenologiche. La criticità di tali sistemi risiede nella necessità di garantire un sincronismo temporale rigoroso tra il campionamento dei sensori e il dato GNSS (\textit{Global Navigation Satellite System})\footnote{Il sistema gestisce la ricezione concorrente di costellazioni GPS (USA), GLONASS (Russia) e Galileo (EU) per massimizzare la precisione del fix in ambienti ostili.}, preservando l'integrità delle informazioni in memorie non volatili anche in condizioni operative isolate o soggette a instabilità energetica \cite{sam_m8q_ds}.

L'applicazione di tali tecnologie spazia dal monitoraggio di micro-climi agricoli per l'ottimizzazione delle risorse (\textit{Smart Farming}), alla tracciabilità di \textit{asset} logistici in tempo reale, fino alla gestione di nodi sensoriali urbani per la mappatura dell'inquinamento atmosferico e acustico.

\subsection{Motivazioni Progettuali}
I microcontrollori \textit{general-purpose} convenzionali manifestano spesso colli di bottiglia prestazionali quando chiamati a gestire simultaneamente bus di comunicazione seriale (UART, I2C, SPI) e task computazionali intensivi. Questa saturazione delle risorse introduce \textit{jitter}\footnote{Variazione statistica indesiderata nel ritardo di propagazione di un segnale periodico rispetto a un riferimento temporale ideale.} indesiderato nei segnali di I/O, degradando la precisione del sistema. 

L'adozione del SoC \textit{Raspberry Pi RP2040} \cite{rp2040_ds} risponde alla necessità di superare tali limitazioni attraverso un'architettura \textit{dual-core} simmetrica\footnote{L'architettura \textit{ARM Cortex-M0+} del RP2040 permette di dedicare un core al processamento dei dati sensoriali e l'altro alla gestione dello stack di comunicazione o del \textit{filesystem}, riducendo le latenze di \textit{interrupt}.} e, soprattutto, mediante l'uso delle unità \textit{Programmable I/O} (PIO) \cite{rp2040_paper}. Queste ultime permettono un approccio di \textit{hardware-software co-design}, dove le macchine a stati dedicate gestiscono il \textit{timing} delle periferiche in modo deterministico, svincolando i core principali dall'onere della gestione dei protocolli a basso livello.



\subsection{Obiettivi del Lavoro}
Il presente elaborato descrive l'intero processo di prototipazione della \textit{scheda}, analizzandone ogni fase critica: come la validazione del \textit{filesystem fail-safe LittleFS}\footnote{\textit{LittleFS} è un \textit{filesystem} progettato specificamente per microcontrollori, dotato di meccanismi di \textit{wear-leveling} e resilienza ai crash energetici (\textit{copy-on-write}).} per la prevenzione della corruzione dei dati.

L'obiettivo primario è la realizzazione di una piattaforma embedded, ottimizzata per il \textit{data-logging} geo-referenziato ad alta frequenza. La documentazione di riferimento, le API sviluppate e i dettagli tecnici sull'architettura firmware sono resi disponibili pubblicamente \cite{ppse_api, ppse_homepage, ppse_github}, unitamente ai report accademici e tecnici di validazione \cite{ppse_academic_report, ppse_technical_ref}. 



Le funzionalità principali del sistema comprendono l'acquisizione di
dati GPS (modulo SAM-M8Q), la sincronizzazione dei relativi dati meteo da un PC host tramite seriale, la lettura di sensori ambientali via
I2C, l'interazione tramite display OLED SSD1306, e garantire la persistenza
su flash. Viene inoltre previsto un feedback visivo su strip WS2812B e acustico tramite buzzer PWM.
